\section{Ermittlung der Alterungsrückstellung}\label{sec:ar}

Dieser Abschnitt regelt die tarifliche Alterungsrückstellung, die
Bilanzdeckungsrückstellung, die Zusatzrückstellungen sowie den
Übertragungswert.

\subsection{Tarifliche Alterungsrückstellung}\label{sec:tarifar}

Die Alterungsrückstellung wird prospektiv als Differenz des Barwerts
der künftigen Leistungen und des Barwerts der künftigen Beiträge
ermittelt (\S\,341f HGB, \S\,3 KVAV; derselbe Rechnungszins, dieselbe
Ausscheideordnung und dieselben Kopfschäden wie für die
Beitragsberechnung). \textbf{Festlegung:} Die Zuschläge $\Omega_x$ und
$\gamma_x$ werden in der Rückstellungsbewertung mit der Stufe des
\emph{erreichten} Alters angesetzt (Stufenwechsel zum Alter~65
entsprechend dem tatsächlichen Wegfall der kostenbedingten Zuschläge),
während die Beitragsformel~\eqref{eq:brutto} die Stufe des
Eintrittsalters verwendet; der hieraus resultierende Sprung der
Rückstellung beim Übergang vom Alter 64 zum Alter 65 wird im Beispiel
des Anhangs~5 ausgewiesen. Für eine versicherte Person mit
erreichtem Alter $x$ und monatlichem Bruttobeitrag $b$ (ohne
gesetzlichen Beitragszuschlag und ohne Risikozuschläge) beträgt die
(gezillmerte) Alterungsrückstellung
\begin{equation}
V_x(b) = G\cdot A_x
- \bigl[\,12\cdot\bigl(1-(\sigma+\Omega_x)\bigr)\cdot b - \gamma_x\,\bigr]
\cdot a_x,
\label{eq:ar}
\end{equation}
wobei
\begin{itemize}\itemsep1pt
  \item $G\cdot A_x$ den Barwert der künftigen Leistungen (EUR),
  \item $12\,(1-(\sigma+\Omega_x))\,b - \gamma_x$ den Jahresbetrag des
        Beitrags nach Abzug der laufenden Zuschläge
        (\glqq Zillmerprämie\grqq),
  \item $a_x$ den Rentenbarwert sowie $\Omega_x$ und $\gamma_x$ die
        Zuschläge zum erreichten Alter
\end{itemize}
bezeichnen. Für den Neuzugang mit Eintrittsalter $x_e$ und
Tarifbeitrag $b = b_{x_e}$ gilt unmittelbar nach
Gleichung~\eqref{eq:brutto}
$V_{x_e}(b_{x_e}) = -\,\alpha_{x_e}\cdot b_{x_e}$ (negativer
Anfangswert in Höhe der Zillmerkosten). Kinder und Jugendliche bilden
keine Alterungsrückstellung. Einmalbeträge (Limitierung, Gutschriften)
erhöhen die Alterungsrückstellung unmittelbar. Ein Beispiel für die
Entwicklung der Alterungsrückstellung enthält Anhang~\ref{sec:anhang5}.

Die jährliche Fortschreibung folgt aus Gleichung~\eqref{eq:ar} als
Rekursion
\begin{equation}
V_{x+1}(b) = \frac{P^{Z}_x - G\,k_x + V_x(b)}{v\cdot\bigl(1-(q+w)_x\bigr)},
\qquad P^{Z}_x = 12\cdot\bigl(1-(\sigma+\Omega_x)\bigr)\cdot b - \gamma_x,
\label{eq:arrek}
\end{equation}
wobei $P^{Z}_x$ die Zillmerprämie (Jahresbetrag des Beitrags nach Abzug
der laufenden Zuschläge), $G\,k_x$ den rechnungsmäßigen Kopfschaden und
der Nenner $v\,(1-(q+w)_x) = DQ_x$ die versicherungsmathematische
Aufzinsung einschließlich der Vererbung der Rückstellungen der durch
Tod und Storno Ausscheidenden bezeichnet (Sterblichkeits- und
Stornogewinne verbleiben rechnungsmäßig im Kollektiv). Die Rekursion
gilt innerhalb konstanter Zuschlagsstufen; am Übergang vom Alter 64
zum Alter 65 ist die Rückstellung nach Gleichung~\eqref{eq:ar} mit den
Zuschlägen des erreichten Alters neu anzusetzen (siehe die obige
Festlegung).

\subsection{Bilanzdeckungsrückstellung}\label{sec:bilanzar}

Für den Bilanzstichtag 31.12. wird die Alterungsrückstellung als
arithmetisches Mittel der Werte zum letzten und nächsten Jahrestag
angesetzt (\S\,18 KVAV; Gleichverteilung der Jahrestage unterstellt).
Hierzu werden die gemittelten Barwerte
\begin{equation}
A_{x+1/2} = \tfrac{1}{2}\bigl(A_x + A_{x+1}\bigr),
\qquad
a_{x+1/2} = \tfrac{1}{2}\bigl(a_x + a_{x+1}\bigr),
\qquad A_{111} := 0,\; a_{111} := 0,
\label{eq:halbbarwerte}
\end{equation}
verwendet (Anlage~\ref{sec:anlage521}); die
Bilanzdeckungsrückstellung einer versicherten Person ergibt sich aus
Gleichung~\eqref{eq:ar} mit $A_{x+1/2}$ und $a_{x+1/2}$ anstelle von
$A_x$ und $a_x$. Ergibt die Aufrechnung negativer gegen positive
Rückstellungen über den Bestand einen negativen Wert, wird in der
Bilanz Null angesetzt (\S\,25 Abs.\,5 RechVersV).

\textbf{Nachweis nach \S\,8 Abs.\,3 KVAV:} Die gezillmerte
Alterungsrückstellung je Eintrittsalter ist spätestens ab dem
\MaxNegM.~Versicherungsjahr (Männer) bzw. \MaxNegF.~Versicherungsjahr
(Frauen) nichtnegativ (Anlage~\ref{sec:anlage522});
damit ist die Gesamtalterungsrückstellung eines Zugangsjahrgangs
höchstens vier Jahre und jede Einzelalterungsrückstellung weder länger
als fünfzehn Jahre noch länger als die Hälfte der tariflich
vorgesehenen künftigen Vertragsdauer negativ.

\subsection{Zusatzrückstellung aus dem gesetzlichen Beitragszuschlag}
\label{sec:zusatzrueckstellung}

Der gesetzliche Beitragszuschlag (Abschnitt~\ref{sec:gz}) wird
individuell in einer Zusatzrückstellung angesammelt (\S\,149, \S\,150
Abs.\,1 VAG). Da künftige Entnahmen nicht statistisch feststehen, wird
die Zusatzrückstellung retrospektiv fortgeschrieben:
\begin{equation}
V^{GZ}_{m+1}
= \frac{\bigl(V^{GZ}_{m} + C_m - E_m\bigr)\cdot(1+i)}{1-q_{x_e+m}},
\qquad V^{GZ}_{0} = 0,
\label{eq:gz-rueckstellung}
\end{equation}
wobei $C_m = 0{,}1\cdot 12\,b$ den Jahresbetrag des Zuschlags
(für $22 \le x_e+m \le 60$, sonst 0; Abschnitt~\ref{sec:gz}), $E_m$
die Entnahmen zur Beitragsermäßigung (ab Alter 65; vorher 0) und
$q_{x_e+m}$ die rechnungsmäßige Sterbewahrscheinlichkeit bezeichnen.
\textbf{Festlegung:} Die Vererbung erfolgt vereinfachend nur über die
Sterblichkeit; bei Storno wird die Zusatzrückstellung -- mangels
Übertragungswertanspruchs in der alten Welt -- zugunsten des
verbleibenden Bestandes vererbt und hierfür im Rahmen der jährlichen
Überschussanalyse gesondert ausgewiesen (eine tafelmäßige Aufteilung
des Stornos nach Zielsystem GKV/PKV liegt nicht vor; der Ansatz ohne
Storno-Vererbung in Gleichung~\eqref{eq:gz-rueckstellung} ist
vorsichtig).

\textbf{Weitere Zusatzrückstellungen nach \S\,150 VAG:}
\begin{itemize}\itemsep2pt
  \item \emph{Zuschreibung von Zinserträgen (\S\,150 Abs.\,1 und 2
        VAG):} Der Rückstellung zur Beitragsermäßigung im Alter werden
        jährlich 90\,\% des rechnungsmäßig ermittelten Überzinses
        gutgeschrieben:
        \begin{equation}
        Z_t = 0{,}9\cdot(r_t - i)\cdot \bar{V}_t,
        \label{eq:ueberzins}
        \end{equation}
        wobei $r_t$ die maßgebliche Kapitalanlagerendite des
        Geschäftsjahres $t$, $i$ den Rechnungszins und $\bar{V}_t$ die
        maßgebliche mittlere positive Alterungsrückstellung des
        Vertrags bezeichnen. Für Versicherte, die das 65.~Lebensjahr
        vollendet haben, werden die Beträge nach Maßgabe des \S\,150
        Abs.\,2 und 3 VAG zur Beitragsermäßigung eingesetzt bzw.
        (über 80-Jährige) jährlich ausgekehrt
        (Abschnitt~\ref{sec:beitragsermaessigung}); der Auskehrbetrag
        entspricht der Zuschreibung $Z_t$ des Geschäftsjahres.
  \item \emph{Rückstellung für erfolgsunabhängige Beitragsrückerstattung
        (\S\,150 Abs.\,4 VAG):} Ihr werden jährlich die vereinnahmten
        euBR-Zuschläge zugeführt und die ausgezahlten
        erfolgsunabhängigen Beitragsrückerstattungen entnommen:
        $\Delta R^{\mathrm{euBR}}_t = \sum \gamma^{\mathrm{euBR}}
        - \text{Auszahlungen}_t$ (Summation über die beitragspflichtigen
        Versicherten; Nachweis der Angemessenheit der
        $\gamma^{\mathrm{euBR}}$-Sätze in
        Abschnitt~\ref{sec:eubr}).
\end{itemize}

\subsection{Übertragungswert}\label{sec:uew}

Der Tarif gehört zur alten Welt (Vertragsabschlüsse vor dem
01.01.2009). Ein Anspruch auf Mitgabe eines Übertragungswertes nach
\S\,146 Abs.\,1 Nr.\,5 VAG i.\,V.\,m. \S\,14 KVAV besteht nicht;
bei einem Wechsel in einen Tarif der neuen Welt beim selben
Unternehmen wird die Alterungsrückstellung nach
Abschnitt~\ref{sec:umstufung} angerechnet und der ab dem Wechsel
aufgebaute Teil gemäß \S\,14 Abs.\,4 KVAV gesondert geführt
(Dokumentation in der TB des Zieltarifs).
